\section{革命时期的地方公众、教育与家庭时间}
咨议局、铁路、军队和退位诏书留下了密集的制度文字，却未自动记录家庭如何获知、承担或回应政治转换。本节在账本的范围内区分官方命令、地方抗议、报刊与电报传播、军队行动和谈判文本，避免用一个中心的时间取代各地社会的不同节奏。
\subsection{1900—1909：地方材料所能看见的尺度}
以下事件按账本时间排列；它们不是同质的社会样本，而是提示应到何种地点、机构和材料中继续追问的入口。
\hypertarget{social:EQ-1909-PROVINCIAL-ASSEMBLIES}{}\paragraph{宣统元年（1909年）}各省咨议局开办，地方精英获得新的公开政治场域。从地方治理、身份、家庭与物质生活的角度，这一节点要求追问它在何处改变了资源、道路、组织或日常风险。核心来源为《宣统政纪》宣统元年相关条；宫中朱批奏折、内阁大库题本与《清史稿》相关志传。它首先留下的是官府分类、任命、处置或资源调拨的痕迹；要判断基层是否照办，仍须寻找县档、账簿、契约或后续案件。账本所示的研究方向是宫廷政治、制度运作与中央—地方信息链研究。可核的条件包括预备立宪需要有限度吸纳地方意见。应分层观察的后果是选举资格、代表性和实际权力均受限制。证据界限：以宣统元年（1909年）的同期文书为定位起点；官修记录、奏报或外交文本服务特定机构立场，具体日期、规模、地方执行与对手视角须以独立材料复核。
\hypertarget{social:EQ-1909-RAILWAY-RECOVERY-MOVEMENT}{}\paragraph{宣统元年（1909年）}铁路国有化和路权问题引发各省商绅、公司和官府的争论。从地方治理、身份、家庭与物质生活的角度，这一节点要求追问它在何处改变了资源、道路、组织或日常风险。核心来源为《宣统政纪》宣统元年相关条；户部奏销、盐政、河工或海关档案。它首先留下的是官府分类、任命、处置或资源调拨的痕迹；要判断基层是否照办，仍须寻找县档、账簿、契约或后续案件。账本所示的研究方向是财政账册、市场网络、灾害环境与区域经济研究。可核的条件包括外债、地方集资与中央财政需求相冲突。应分层观察的后果是“收回路权”诉求包含不同地区和利益群体。证据界限：以宣统元年（1909年）的同期文书为定位起点；官修记录、奏报或外交文本服务特定机构立场，具体日期、规模、地方执行与对手视角须以独立材料复核。
\subsection{1910—1919：地方材料所能看见的尺度}
以下事件按账本时间排列；它们不是同质的社会样本，而是提示应到何种地点、机构和材料中继续追问的入口。
\hypertarget{social:EQ-1910-ZIZHENG-YUAN}{}\paragraph{宣统二年（1910年）}资政院开院，中央层面的立宪咨询机构开始运作。从地方治理、身份、家庭与物质生活的角度，这一节点要求追问它在何处改变了资源、道路、组织或日常风险。核心来源为《宣统政纪》宣统二年相关条；宫中朱批奏折、内阁大库题本与《清史稿》相关志传。它首先留下的是官府分类、任命、处置或资源调拨的痕迹；要判断基层是否照办，仍须寻找县档、账簿、契约或后续案件。账本所示的研究方向是宫廷政治、制度运作与中央—地方信息链研究。可核的条件包括朝廷试图展示宪政进程并控制代表参与。应分层观察的后果是咨询权与立法权的界限成为持续争议。证据界限：以宣统二年（1910年）的同期文书为定位起点；官修记录、奏报或外交文本服务特定机构立场，具体日期、规模、地方执行与对手视角须以独立材料复核。
\hypertarget{social:EQ-1910-FAMINE-AND-UNREST}{}\paragraph{宣统二年（1910年）}灾荒、物价和财政压力持续影响晚清地方社会与政治动员。从地方治理、身份、家庭与物质生活的角度，这一节点要求追问它在何处改变了资源、道路、组织或日常风险。核心来源为《宣统政纪》宣统二年相关条；地方志、督抚奏折及地方档案。这类编年材料可以标定朝廷获知和叙述事件的时间，却往往压低妇女、儿童、佃户、流民和非文书精英的行动。账本所示的研究方向是地方档案、迁徙、宗教、族群与社会动员研究。可核的条件包括自然环境、市场波动和行政能力交互作用。应分层观察的后果是地方灾情与政治行动不可简单建立单因果关系。证据界限：以宣统二年（1910年）的同期文书为定位起点；官修记录、奏报或外交文本服务特定机构立场，具体日期、规模、地方执行与对手视角须以独立材料复核。
\hypertarget{social:EQ-1911-RAILWAY-NATIONALIZATION}{}\paragraph{宣统三年（1911年）}清廷宣布铁路干线国有化并以外债筹资，保路运动迅速扩大。从地方治理、身份、家庭与物质生活的角度，这一节点要求追问它在何处改变了资源、道路、组织或日常风险。核心来源为《宣统政纪》宣统三年相关条；宫中朱批奏折、内阁大库题本与《清史稿》相关志传。它首先留下的是官府分类、任命、处置或资源调拨的痕迹；要判断基层是否照办，仍须寻找县档、账簿、契约或后续案件。账本所示的研究方向是宫廷政治、制度运作与中央—地方信息链研究。可核的条件包括中央财政、外债与地方商办铁路利益冲突。应分层观察的后果是政策公告、地方抗议和后续镇压应分开记录。证据界限：以宣统三年（1911年）的同期文书为定位起点；官修记录、奏报或外交文本服务特定机构立场，具体日期、规模、地方执行与对手视角须以独立材料复核。
\hypertarget{social:EQ-1911-SICHUAN-RAILWAY-PROTEST}{}\paragraph{宣统三年（1911年）}四川保路运动升级，赵尔丰处置引发更大政治危机。从地方治理、身份、家庭与物质生活的角度，这一节点要求追问它在何处改变了资源、道路、组织或日常风险。核心来源为《宣统政纪》宣统三年相关条；地方志、督抚奏折及地方档案。这类编年材料可以标定朝廷获知和叙述事件的时间，却往往压低妇女、儿童、佃户、流民和非文书精英的行动。账本所示的研究方向是地方档案、迁徙、宗教、族群与社会动员研究。可核的条件包括地方股东、商绅、学生和官府对路权与财政有不同诉求。应分层观察的后果是参与规模和暴力过程须依当地档案与报刊核验。证据界限：以宣统三年（1911年）的同期文书为定位起点；官修记录、奏报或外交文本服务特定机构立场，具体日期、规模、地方执行与对手视角须以独立材料复核。
\hypertarget{social:EQ-1911-WUCHANG-UPRISING}{}\paragraph{宣统三年（1911年）}武昌起义爆发，新军与革命团体控制武汉部分地区。从地方治理、身份、家庭与物质生活的角度，这一节点要求追问它在何处改变了资源、道路、组织或日常风险。核心来源为《宣统政纪》宣统三年相关条；军机处档、战事方略及地方奏报。这类编年材料可以标定朝廷获知和叙述事件的时间，却往往压低妇女、儿童、佃户、流民和非文书精英的行动。账本所示的研究方向是战争动员、军需、战场暴力与地方社会研究。可核的条件包括新军组织、革命网络和铁路危机相互叠加。应分层观察的后果是起义的即时参与者、城市控制与后续省份响应需分层叙述。证据界限：以宣统三年（1911年）的同期文书为定位起点；官修记录、奏报或外交文本服务特定机构立场，具体日期、规模、地方执行与对手视角须以独立材料复核。
\hypertarget{social:EQ-1911-PROVINCIAL-INDEPENDENCE}{}\paragraph{宣统三年（1911年）}多省宣布脱离清廷，帝国行政和财政网络快速分裂。从地方治理、身份、家庭与物质生活的角度，这一节点要求追问它在何处改变了资源、道路、组织或日常风险。核心来源为《宣统政纪》宣统三年相关条；宫中朱批奏折、内阁大库题本与《清史稿》相关志传。它首先留下的是官府分类、任命、处置或资源调拨的痕迹；要判断基层是否照办，仍须寻找县档、账簿、契约或后续案件。账本所示的研究方向是宫廷政治、制度运作与中央—地方信息链研究。可核的条件包括武昌起义、地方军政力量和立宪派选择共同推动变化。应分层观察的后果是“独立”声明不等于各省形成稳定统一控制。证据界限：以宣统三年（1911年）的同期文书为定位起点；官修记录、奏报或外交文本服务特定机构立场，具体日期、规模、地方执行与对手视角须以独立材料复核。
\hypertarget{social:EQ-1911-YUAN-SHIKAI-RETURN}{}\paragraph{宣统三年（1911年）}清廷起用袁世凯处理北方军政和议和事务。从地方治理、身份、家庭与物质生活的角度，这一节点要求追问它在何处改变了资源、道路、组织或日常风险。核心来源为《宣统政纪》宣统三年相关条；宫中朱批奏折、内阁大库题本与《清史稿》相关志传。它首先留下的是官府分类、任命、处置或资源调拨的痕迹；要判断基层是否照办，仍须寻找县档、账簿、契约或后续案件。账本所示的研究方向是宫廷政治、制度运作与中央—地方信息链研究。可核的条件包括中央缺乏可依赖的军队与谈判人选。应分层观察的后果是袁在清廷、北洋与革命派之间的策略塑造退位进程。证据界限：以宣统三年（1911年）的同期文书为定位起点；官修记录、奏报或外交文本服务特定机构立场，具体日期、规模、地方执行与对手视角须以独立材料复核。
\hypertarget{social:EQ-1912-REPUBLIC-FOUNDING}{}\paragraph{宣统四年（1912年）}中华民国临时政府在南京成立，与清廷退位谈判并行。从地方治理、身份、家庭与物质生活的角度，这一节点要求追问它在何处改变了资源、道路、组织或日常风险。核心来源为《宣统政纪》宣统四年相关条；宫中朱批奏折、内阁大库题本与《清史稿》相关志传。它首先留下的是官府分类、任命、处置或资源调拨的痕迹；要判断基层是否照办，仍须寻找县档、账簿、契约或后续案件。账本所示的研究方向是宫廷政治、制度运作与中央—地方信息链研究。可核的条件包括各省独立与南北议和创造新的政权竞争格局。应分层观察的后果是临时政府的法理主张与实际控制范围需要区分。证据界限：以宣统四年（1912年）的同期文书为定位起点；官修记录、奏报或外交文本服务特定机构立场，具体日期、规模、地方执行与对手视角须以独立材料复核。
\hypertarget{social:EQ-1912-QING-ABDICATION}{}\paragraph{宣统四年（1912年）}宣统帝颁布退位诏书，清帝国终结。从地方治理、身份、家庭与物质生活的角度，这一节点要求追问它在何处改变了资源、道路、组织或日常风险。核心来源为《宣统政纪》宣统四年相关条；宫中朱批奏折、内阁大库题本与《清史稿》相关志传。它首先留下的是官府分类、任命、处置或资源调拨的痕迹；要判断基层是否照办，仍须寻找县档、账簿、契约或后续案件。账本所示的研究方向是宫廷政治、制度运作与中央—地方信息链研究。可核的条件包括南北谈判、财政军政压力与皇室待遇安排共同促成退位。应分层观察的后果是退位文本规定的优待与地方政治现实须分别追踪。证据界限：以宣统四年（1912年）的同期文书为定位起点；官修记录、奏报或外交文本服务特定机构立场，具体日期、规模、地方执行与对手视角须以独立材料复核。
